% !TEX root = ../../cookbook.tex
\newpage
\recipe{Risotto Patrio}{{\spaceskip=0.9\fontdimen2\font Perfected, in my not-so-humble opinion, from \href[]{https://www.gadda.ed.ac.uk/Pages/resources/essays/risotto.php}{Carlo Emilio Gadda's version}. Serves 4.}}{}{}

\begin{multicols}{2}
	\subsection*{Ingredients}
	\begin{ingredients}
		\item 1 onion, finely diced
		\item 350 g risotto rice\footnote{The Vialone type is ideal, but the Arborio is easier to find in the US.}
		\item High-quality Lodi butter, \textit{quantum prodest}
		\item 1.5 L richly flavored beef stock
		\item 0.2 g (or 50 threads) saffron
		\item White wine, for deglazing
		\item Grated Parmigiano Reggiano
		\item Porcini, chanterelles, or morel mushrooms
	\end{ingredients}
	
	\columnbreak
	
	\subsection*{Equipment}
	\begin{equipment}
		\item Tinned copper saucepan (or aluminum, but do not use a non-stick pot)
		\item Ladle
		\item Wooden spoon
	\end{equipment}
	\vfill
\end{multicols}

\medskip

% --- Preparation ---
In a separate bowl or pot, prepare 1.5 L of \textbf{saffron-infused beef stock} so that it is ready to use.
If available, add the \textbf{mushrooms} to the stock as well.
In the saucepan, melt some \textbf{butter}\footnote{Optionally, add a tablespoon of olive oil so that the butter is less likely to burn.} and gently cook the finely diced \textbf{onion}.
As soon as the buttery onion mixture begins to sizzle, add the \textbf{rice} and toast it, stirring against the copper bottom.
Deglaze with the \textbf{white wine}.
Then begin cooking the rice by adding the stock a ladleful at a time until the rice is cooked: not overdone, but just beyond \textit{al dente}.
Remove the risotto from the heat, \textit{mantecate}\footnote{To \textit{mantecare} a risotto means to vigorously stir in cold butter and grated Parmigiano Reggiano off the heat until the risotto becomes smooth, creamy, and glossy.} it, and add as much \textbf{Parmigiano Reggiano} as you like.
Cover with a lid and let it rest for 10 minutes before serving\footnote{Best enjoyed with some rustic bread and a glass of red wine.}. 

% --- Description ---
\begin{recipeblock}
	
\end{recipeblock}
